\section{Related Work}
\label{sec:related}

\subparagraph*{Satisfiability and entailment.}
Subtype constraints over trees were introduced for type inference with
subtyping \cite{Mitchell91,FuhMishra90}. For non-structural subtyping,
satisfiability is decidable in polynomial time over finite and infinite trees
\cite{KozenPalsbergSchwartzbach94,PalsbergWandOKeefe97}; over finite
structural trees it depends on the poset of base types
\cite{Tiuryn92,PrattTiuryn96,Frey97,NiehrenPriesnitzSu05}. Entailment was
studied for constraint simplification \cite{Pottier96,TrifonovSmith96,
Rehof97,Pottier01}. For structural subtyping it is coNP-complete over finite
trees and PSPACE-complete over infinite trees
\cite{HengleinRehof97,HengleinRehof98}. For non-structural subtyping,
Henglein and Rehof proved PSPACE-hardness \cite{HengleinRehof98}; Niehren
and Priesnitz characterized entailment by automata with comparisons between
suffixes of paths, proved PSPACE-completeness for a fragment and decided the
unary case \cite{NiehrenPriesnitz99,NiehrenPriesnitz03}; decidability, TLCA
Problem~16 \cite{Rehof96TLCA}, was proved in \cite{Zhao26}. Existential
entailment, which underlies the subsumption of constrained types
\cite{TrifonovSmith96}, is decided in \cite{Zhao26E}. The present paper uses
the leaf cones of \cite{Zhao26E} for the reduction to automata and follows
\cite{Zhao26} in reducing decidability to a bound on shortest holes; the
readers with several starts, the treatment of rejected tails and the
simultaneous shortening of tuples are new.

\subparagraph*{The first-order theory.}
Su, Aiken, Niehren, Priesnitz and Treinen proved that the first-order theory
of non-structural subtyping is undecidable, for finite and for infinite
trees, and decidable when all constructors are unary or nullary
\cite{SuAikenNiehrenPriesnitzTreinen02,SuEtAl06}. Their sentences have
several quantifier alternations; undecidability already holds for the prefix
$\exists\forall\exists$, for a binary covariant constructor with or without
$\Top$, over finite, regular and arbitrary trees \cite{Zhao26U}. For structural subtyping of
non-recursive types the full first-order theory is decidable
\cite{KuncakRinard03}. For non-structural subtyping with one binary
covariant constructor, our result places the boundary: the $\Sigma_1$ level
is decidable, the $\Sigma_3$ level is not, and the level in between is open
except for the existential entailments of \cite{Zhao26E}. For unary and
nullary constructors the whole theory is decidable. Su \cite{Su02} reduces
first-order questions to the emptiness of constrained tree automata, which
is undecidable in general; our readers are a decidable special case in which
the tests compare rests of words.

\subparagraph*{Ordering constraints over feature trees.}
For feature trees ordered by information, entailment and the existential
fragment of the first-order theory are well understood: entailment with
existential quantifiers is PSPACE-complete, and the full first-order theory
is undecidable \cite{MullerNiehrenTreinen98,MullerNiehrenTreinen01,
MullerNiehrenPodelski00}. Feature trees have an open alphabet of features,
so a countermodel can always use a fresh feature. A ranked signature offers
no fresh direction, and this is the source of the comparisons between
different witness paths in our holes.

\subparagraph*{Word equations and combinatorics on words.}
Niehren and Priesnitz observed that entailment involves word equations of a
restricted kind \cite{NiehrenPriesnitz03}. The existential theory of words
with concatenation is decidable by Makanin's algorithm
\cite{Makanin77,Plandowski04,Schulz92}, but the conditions of a hole
quantify universally over the cuts of the unknown words, and theories of
words with a quantifier alternation are undecidable \cite{Durnev95}.
Our proof uses elementary combinatorics on words:
the periodicity lemma of Fine and Wilf \cite{FineWilf65,Lothaire02}, squares
and maximal repetitions \cite{CrochemoreRytter95,KolpakovKucherov99}, and an automaton
of suffixes in the style of Aho and Corasick \cite{AhoCorasick75}. The
marker hierarchy of \cref{sec:fresh} is a parsing of several words by their
own suffixes, in the spirit of the hierarchical decompositions used for
plain entailment \cite{Zhao26}.

\subparagraph*{Other approaches to subtyping.}
Algebraic subtyping \cite{DolanMycroft17,Dolan16,ParreauxChau22} changes the
model of types so that constraints can be solved and types simplified by
automata-theoretic means, and semantic subtyping
\cite{CastagnaXu11,GesbertGenevesLayaida15} interprets types as sets, with
decidable subtyping in the presence of Boolean connectives. These systems
have union, intersection and often negation \emph{types}; our sentences have
disjunction and negation of \emph{constraints} over the classical tree model
with $\Bot$, $\Top$ and one constructor.
